% !PRINTED 396
\weberpage{396}
12. \gesperrt{We divide the ideals $\mathfrak{a}$ into genera $G, G_1, G_2, \ldots$ by uniting all ideals $\mathfrak{a}_{\nu}$ whose norms $N(\mathfrak{a}_{\nu})$ are contained in one class $R_{\nu}$ into a genus $G_{\nu}$.}

13. \gesperrt{The ideals whose norms are at the same time norm residues $r$ belong to the principal genus $G$, and the principal genus is itself a group.}

If we wish to take the orders into account at the same time, we must also admit equivalence with respect to the order $[Q]$ (\S\,100), which of course (for $Q = 1$) includes general equivalence. Then the following holds:

14. \gesperrt{The ideals of the principal class belong to the principal genus, and equivalent ideals have the same genus.}

For if $\mathfrak{a}$ is equivalent to $\mathfrak{a}_1$, then there is a number $\eta$ in $[Q]$ for which $\mathfrak{a} = \eta \mathfrak{a}_1$; consequently $N(\mathfrak{a}) = N(\eta) N(\mathfrak{a}_1)$, and $N(\eta)$ is a number in $R$. Consequently $\mathfrak{a}$ and $\mathfrak{a}_1$ belong, by 13., to the same genus. The genera therefore comprise not merely the individual ideals, but the \gesperrt{ideal classes}.

In order that a number $a$ from $Z$ can be an ideal norm, a certain condition must be fulfilled. For if
\weq{3}{a = N(\mathfrak{a}),}
\gesperrt{and $\mathfrak{a}$ is in the first place a primary integral ideal, then by \S\,91}
\weq{4}{\Delta = b^2 - 4 a c,}
where $b, c$ are integral rational numbers, and consequently
\weq{5}{(\Delta, a) = +1.}
If $\mathfrak{a}$ is not primary, but equal to $m \mathfrak{a}_0$ with primary $\mathfrak{a}_0$, and $a_0 = N(\mathfrak{a}_0)$, then $a = m^2 a_0$, and consequently, by \S\,85, (15), condition (5) must be satisfied here also; and the same follows for fractional numbers as well, if numerator and denominator are relatively prime to $D$, provided one regards $a$ in (5) as an integral representative of its class with respect to the modulus $D$.

Since there are now numbers $a$ for which $(\Delta, a) = -1$, not every number in $Z$ can be the norm of an ideal, and the number of the genera is therefore smaller than $\mu$. It must be a divisor of $\mu$, i.e. a power of 2, and consequently we have the theorem:

15. \gesperrt{The number $g$ of the genera is at most equal to}
\[
\frac{1}{2} (Z, R),
\]
\gesperrt{that is, to half of the number determined in 9.}
